% Mahican (Edwards 1787) — Source Workflow
% Issue #1369
% Status: Research complete, implementation pending

\section{Mahican: Edwards 1787}
\label{sec:1369}

\subsection{Source Description}

\subsubsection{Jonathan Edwards Jr.\ --- Biographical Context}

Jonathan Edwards Jr.\ (May 26, 1745 -- August 1, 1801) was a
Congregationalist minister, theologian, college president, and linguist. He
was the son of the theologian Jonathan Edwards Sr.\ (Princeton's 3rd
president). He graduated from Princeton University in 1765 and served as
president of Union College in Schenectady, New York from 1799 until his death
in office.

Edwards grew up in Stockbridge, Massachusetts, a ``praying town'' where
Mahican people were the majority population (approximately 150 Mahican
families versus approximately 12 white families). He learned Mahican as a
child from playmates, becoming fluent. At age 10 (1755), his father sent him
to live in the Iroquois settlement of Onaquaga (Onohoquaga) to learn Oneida,
giving him first-hand knowledge of both Algonquian and Iroquoian language
families.

Unlike his father, who refused to learn Mahican (calling it ``barbarous''),
Edwards Jr.\ defended the language's complexity, demonstrated that it had
abstract terms (love, hatred, malice, religion), and translated the Lord's
Prayer into Mahican. He consulted a Mahican elder, ``Captain Yohum,'' for
revisions before publication. George Washington, James Madison, and Thomas
Jefferson all received copies of his treatise.

\subsubsection{The 1788 Publication}

\textbf{Full title:} \emph{Observations on the Language of the Muhhekaneew
Indians; in which the Extent of that Language in North-America is shewn; its
Genius Grammatically Traced; some of its Peculiarities, and some Instances
of Analogy between that and the Hebrew are Pointed Out.}

Communicated to the Connecticut Society of Arts and Sciences on October 23,
1787. Published 1788 by Josiah Meigs, New Haven. A second edition with notes
by John Pickering appeared in 1823 (in Massachusetts Historical Collections).

The treatise includes:

\begin{itemize}
  \item A grammatical sketch of Mahican (parts of speech, person marking,
    tense)
  \item A comparative glossary: English--Mahican--Shawnee
  \item A comparative glossary: English--Mahican--Chippeway/Ojibwa
  \item Numerals 1--10 in Mahican and Mohawk
  \item The Lord's Prayer in Mahican and in the ``language of the Six
    Nations''
  \item Word lists organized by semantic domain
  \item Observations on the geographic extent of Algonquian languages
  \item Speculative comparison with Hebrew (a common 18th-century trope)
\end{itemize}

Edwards uses the term ``Mohegan'' interchangeably with ``Muhhekaneew'' for
the language, though he distinguishes it clearly from Mohegan-Pequot
(Southern New England). His dialect is specifically Stockbridge Mahican.

\subsubsection{Orthography System}

Edwards uses an English-based ad-hoc orthography typical of 18th-century
colonial sources. He provides a pronunciation key indicating his intended
sound values for vowel letters.

\textbf{Consonant values (relatively clear):}

\begin{table}[h]
\centering
\begin{tabular}{lll}
\toprule
\textbf{Grapheme} & \textbf{Likely value} & \textbf{Evidence} \\
\midrule
\textit{p, t, k} & /p, t, k/ (unaspirated or lightly aspirated) & PA *p, *t, *k $>$ Mahican /p, t, k/ \\
\textit{b, d, g} & Probably /p, t, k/ after nasals or in voiced contexts & Unclear; English-orthography convention \\
\textit{m, n} & /m, n/ & Straightforward \\
\textit{s} & /s/ & PA *s, *š merged to /s/ in Mahican by modern period \\
\textit{sh} & /š/ & English convention \\
\textit{h} & /h/ & Explicitly written; Edwards is unusually consistent here \\
\textit{w, y} & /w, j/ & Glides \\
\textit{ch} & /č/ & English convention \\
\textit{gh} & \textbf{UNKNOWN} & See \S\ref{sec:1369-gh} \\
\bottomrule
\end{tabular}
\caption{Edwards' consonant graphemes and likely values}
\label{tab:1369-consonants}
\end{table}

\textbf{Vowel values (more ambiguous):}

The key insight (Goddard 2008, Masthay 1991, Warne) is that Edwards uses
\textit{e} to represent /iː/ in many cases. This is confirmed by his
pronunciation key and by cognate comparison.

\begin{table}[h]
\centering
\begin{tabular}{llll}
\toprule
\textbf{Edwards grapheme} & \textbf{Likely Mahican phoneme} & \textbf{PA source} & \textbf{Example} \\
\midrule
\textit{ee} & /iː/ & PA *iː & \textit{neesoh} `two' (PA *niːšwi) \\
\textit{e} & /iː/ or /i/ & PA *i(ː) & \textit{seepoo} `river' (PA *siːpoːwi) \\
\textit{a} & /a/ or /aː/ & PA *a, *eː $\rightarrow$ PEA *aː & \textit{nannah} `five' (PA *nyaːlanwi) \\
\textit{o, oo} & /o/ or /u/ & PA *o(ː) & \\
\textit{u} & /ə/ (schwa) in unstressed; possibly /o/ elsewhere & PA *e $\rightarrow$ PEA *ə & \textit{muhhekaneew} (with mə h-) \\
\textit{au, aw} & /a/ + /w/ or /aː/ & & \\
\textit{uh} & /ə h/ or /ə/ & & Frequent in final syllables \\
\bottomrule
\end{tabular}
\caption{Edwards' vowel graphemes and likely values}
\label{tab:1369-vowels}
\end{table}

\subsection{Phonemic Reconstruction}

\subsubsection{Key Phonological Correspondences (PA $\rightarrow$ Mahican)}

From Warne's analysis of the Schmick manuscript, confirmed across sources:

\begin{table}[h]
\centering
\begin{tabular}{llll}
\toprule
\textbf{PA} & \textbf{Old Mahican} & \textbf{Modern Mahican} & \textbf{Notes} \\
\midrule
*p & /p/ & /p/ & No change \\
*t & /t/ & /t/ & No change \\
*k & /k/ & /k/ & No change \\
*č & /č/ & /č/ & No change \\
*θ{} & /n/ & /n/ & \textbf{Merged} with *n by Old Mahican period \\
*l & /n/ & /n/ & \textbf{Merged} with *n by Old Mahican period \\
*n & /n/ & /n/ & \\
*s & /s/ & /s/ & \\
*š & /s/ & /s/ & Merged to /s/ by Modern Mahican \\
*m & /m/ & /m/ & \\
*h & /h/ & /h/ & \\
*iː & /iː/ & /iː/ & \\
*i & /i/ & /i/ & \\
*eː & /a/ (lax low) & /a/ & \textbf{Retracted/lowered} from *eː to /a/ in PEA \\
*e & /ə/ (schwa) & /ə/ & Reduced to schwa (pan-EA development) \\
*aː & /aː\~ãː/ & /ãː/ & Nasalized in modern Mahican \\
*a & /a/ & /a/ & \\
*oː & /uː/ & /uː/ & \\
*o & /o/ & /o/ & Marginal phoneme even in PA \\
\bottomrule
\end{tabular}
\caption{PA $\rightarrow$ Mahican sound correspondences}
\label{tab:1369-correspondences}
\end{table}

\subsubsection{The /iː/ vs.\ /eː/ Issue}

Recent work (documented in \emph{ImportantMahicanVowel2.pdf} from
munseedelaware.com, with consultation by Goddard and Masthay) has established
that Edwards' \textit{e} frequently represents /iː/, NOT /eː/. Goddard
confirmed (c.\ 2023 email, cited in the Mahican vowel paper) that the second
vowel of the tribal name was absolutely not /eː/ and is properly the ``lower
high'' version of the vowel /iː/, produced by setting the throat and tongue
in a position to say the English word ``hit'' and then producing the sound.

Evidence for this mapping:

\begin{itemize}
  \item Edwards writes \textit{Muhhekaneew}---the second vowel is /iː/
  \item \textit{Neesoh} `two'---cf.\ PA *niːšwi, Munsee /niːšə/
  \item \textit{Seepoo} `river'---cf.\ PA *siːpoːwi, Munsee /siːpə w/
  \item Goddard's 2008 paper transcribes Dick's /kīsōx/ `sun' and /θ{}īpə w/
    `river' with /iː/---even where the English spelling suggests ``e''
  \item Goddard's 2010 phoneme inventory for Mahican lists NO /eː/ phoneme
\end{itemize}

\textbf{Impact on orthography conversion:} Edwards' \textit{e} should NOT be
mapped to /e/ or /ɛ/ by default. The default mapping should be /i(ː)/ unless
cognate evidence specifically indicates a different value.

\subsection{The \textit{gh} Grapheme Mystery}
\label{sec:1369-gh}

Edwards uses \textit{gh} in several words where the phonological value is
uncertain. Cross-linguistic comparison suggests multiple possible values:

\begin{table}[h]
\centering
\begin{tabular}{lll}
\toprule
\textbf{Candidate} & \textbf{Evidence For} & \textbf{Evidence Against} \\
\midrule
/x/ (voiceless velar fricative) & Schmick uses German \textit{ch} for /x/ in cognate forms; PA *h before consonant sometimes appears as /x/ in daughter languages & \\
/h/ & Simple /h/ is the default reflex of PA *h & Edwards also writes \textit{h} separately, so \textit{gh} may be a distinct phoneme \\
/k/ (lenis/voiced allophone) & English \textit{gh} occasionally used for /k/ in colonial sources & Unlikely given that Edwards has \textit{k} for /k/ \\
\bottomrule
\end{tabular}
\caption{Candidate values for the \textit{gh} grapheme}
\label{tab:1369-gh}
\end{table}

\textbf{Recommended resolution:} Use comparative evidence. If a cognate from
Schmick, Munsee, or Wampanoag shows /x/, the mapping is /x/. If it shows /h/,
the mapping is /h/. If no cognate exists, flag as ambiguous.

\subsection{Phonological Processes Affecting Edwards' Orthography}

\subsubsection{Vowel Reduction and \textit{uh}}

Edwards frequently writes \textit{u} and \textit{uh} in unstressed syllables.
The PA *e regularly became schwa /ə/ in Proto-Eastern Algonquian. In
Edwards' orthography, unstressed \textit{u} most likely = /ə/;
\textit{uh} may represent /ə h/ or breathy schwa in a checked syllable.
Munsee shows parallel reduction: PA *e $\rightarrow$ PEA *ə
$\rightarrow$ Munsee /ə/ (raised to /i/ or /u/ before certain clusters).

\subsubsection{Nasalization of *aː}

PA *aː developed a nasal quality in Mahican, marked by Schmick with a tilde
(ã). Warne considers this a ``secondary feature'' that was optional in Old
Mahican but generalized to ALL instances of /aː/ by the Modern Mahican
period. Edwards does NOT mark nasalization, so Edwards' \textit{a} may
represent either oral /aː/ or nasal /ãː/ depending on cognate evidence.

\subsubsection{Word-Final Vowel Deletion}

PA word-final vowels were deleted in Mahican (a pan-Eastern Algonquian
process). Edwards sometimes writes these (reflecting an older/more careful
pronunciation) and sometimes omits them. For example, PA *-i (plural marker)
sometimes appears as \textit{-eh} or \textit{-uh} in Edwards, sometimes
absent.

\subsubsection{Schwa Lowering Before /h, x/}

Before /h/ and /x/, schwa /ə/ is lowered to /a/ in Mahican (Warne).
This accounts for alternations where Edwards writes \textit{a} where PA
reconstruction predicts *e. Example: \textit{ntah} `my heart' $<$ PA *-teːh-
(with *eː $\rightarrow$ *a and schwa-lowering in certain environments).

\subsubsection{Rounding Assimilation}

PA *e rounds to Mahican /o/ when adjacent to *w. This explains forms like
the word for `tree' and `beaver' where Edwards' \textit{o} corresponds to PA
*e. (Warne's rule ordering: rounding before laxing.)

\subsection{Morphological Patterns in Edwards' Data}

Edwards' data reveals several regular morphological patterns that can aid
word segmentation before orthography conversion.

\subsubsection{Person Prefixes}

\begin{table}[h]
\centering
\begin{tabular}{llll}
\toprule
\textbf{Person} & \textbf{Prefix} & \textbf{Edwards example} & \textbf{Gloss} \\
\midrule
1st sg. & n- & \textit{ncmeetseh, ncmeetfeh} & `I eat' \\
2nd sg. & k- & \textit{kcmeetseh, ka meetjeh} & `you eat' \\
3rd sg. & w-/∅- & \textit{pumiffoo} & `he walks' \\
\bottomrule
\end{tabular}
\caption{Person prefixes in Edwards' data}
\label{tab:1369-prefixes}
\end{table}

Edwards himself noted these correspondences and compared them to Hebrew
prefixes, which was one of his major observations.

\subsubsection{Verbal Suffixes}

Edwards records frequent verbal suffixes: \textit{-uh, -eh, -seh} (likely
aspect/mood markers), \textit{-oo} (3rd person singular independent
indicative, e.g., \textit{pumiffoo}), and \textit{-nuh} (1st person plural,
e.g., \textit{ncmeetfchnuh} `we eat').

\subsubsection{Noun Inflection}

Possessive prefixes precede nouns (e.g., \textit{noghkuh} `my father').
Plural formation occurs via suffixation (e.g., \textit{penumpaufoo} `boy'
$\rightarrow$ \textit{penumpaufoouk} `boys').

\subsection{Comparative Sources for Orthography Interpretation}

\subsubsection{The Schmick Manuscript (Primary Internal Source)}

\begin{itemize}
  \item \textbf{Author:} Johann Jacob Schmick, Moravian missionary
  \item \textbf{Period:} c.\ 1753--1767 (before Edwards)
  \item \textbf{Dialect:} Moravian/Western Mahican (Ohio/PA branch)
  \item \textbf{Format:} Two volumes (322 pp.), German-based orthography
  \item \textbf{Publication:} Masthay (ed.), \emph{Schmick's Mahican
    Dictionary}, APS Memoir 197 (1991)
  \item \textbf{Contains:} Mahican historical phonology by David Pentland
\end{itemize}

Schmick's manuscript uses a German-based orthography which is in some ways
more precise than Edwards' English-based system: German \textit{ch} for /x/
(vs.\ Edwards' ambiguous \textit{gh}), more consistent vowel representation,
and extensive vocabulary (much larger than Edwards' 144 records). However,
Schmick shows inconsistency with vowels---especially the reflex of PA *aː
which he writes variously as \textit{a, aa, ã, ã, a} (Warne). The tilde
likely indicates a secondary nasalization feature.

\textbf{Key finding:} Schmick is much more consistent with consonants than
vowels. Identical forms appear written with different vowel symbols or
sequences (Warne, \emph{Time-Depth in Mahican Diachronic Phonology}).

Pentland's historical phonology (in Masthay 1991) provides the connecting
framework between PA reconstructions and Mahican forms. Goddard (2007--2008)
provided important corrections to Pentland's phonology, particularly in the
treatment of vowel length and the PA *e reflexes.

\subsubsection{Zeisberger's Delaware Materials}

David Zeisberger (1721--1808), Moravian missionary, produced extensive
documentation of Munsee Delaware (a close relative of Mahican within the
Delawaran subgroup): \emph{Zeisberger's Delaware Vocabulary of 1772} (earliest
extensive Munsee vocabulary), \emph{Grammar of the Language of the Lenni
Lenape or Delaware Indians} (1827), and the \emph{Diary of David Zeisberger}
(several volumes). These are important for comparative cognate work because
Munsee and Mahican share a common ancestor (Delawaran, per Goddard) and often
show parallel developments.

\subsubsection{Later Mahican Sources}

\begin{itemize}
  \item \textbf{Hendrick Aupaumut writings} (c.\ 1790s--1820s): Stockbridge
    Mahican letters and documents by a fluent Mahican speaker
  \item \textbf{Wm.\ Dick (Michelson notes)} (c.\ 1910s): Modern Mahican
    (Wisconsin) phonetic transcriptions
  \item \textbf{Bernice Robinson recordings} (mid-20th century): Audio
    recordings of the last speakers
\end{itemize}

\subsection{Problems Encountered}

\begin{itemize}
  \item \textbf{The \textit{gh} grapheme mystery:} \textit{ghusooh} (eight)
    contains \textit{gh}---the most unusual grapheme across all eight sources.
    Candidates include /x/ (voiceless velar fricative), /h/, or /k/. Must be
    resolved per cognate on a case-by-case basis. This is the single biggest
    ambiguity.
  \item \textbf{\textit{uh} reduction:} Edwards' orthography is dominated by
    \textit{u} and \textit{uh} sequences representing reduced /ə/. May
    represent /ə/, /ə h/, or /uh/ depending on context.
  \item \textbf{Word-initial consonant clusters:} Highly unusual by English
    standards (\textit{wnissoo}, \textit{hkeesque}, \textit{hpoon},
    \textit{kmeetseh})---indicate rich Mahican morphology. Need morphological
    analysis to determine whether initial CC sequences are prefix+root or
    root-internal.
  \item \textbf{Sentence-level entries:} Many entries are full sentences, not
    lexemes---require morphological segmentation before conversion.
  \item \textbf{Vowel quality ambiguity:} Edwards' English-based vowel letters
    are ambiguous. The /iː/ = \textit{e} mapping (Goddard) is the most
    important correction. \textit{a} could be /a/, /aː/, or /ãː/.
  \item \textbf{No phonetic transcriptions:} Edwards did not use a systematic
    phonetic alphabet. No phonetic entries exist.
  \item \textbf{Small corpus:} Approximately 144 records limits statistical
    approaches.
  \item \textbf{Dialect differences:} Stockbridge Mahican (Edwards) vs.\
    Moravian Mahican (Schmick) may have phonological differences.
  \item \textbf{18th-century conventions:} Edwards' spellings reflect his own
    English pronunciation as much as Mahican sounds; e.g., \textit{oo} may be
    /uː/ or /oː/.
\end{itemize}

\subsection{Research Approach}

The conversion pipeline follows this sequence:

\begin{enumerate}
  \item \textbf{Preprocessing:} Split multi-word entries into individual
    tokens. Flag entries with \textit{gh} for comparative cognate resolution.
    Normalize Edwards' typographic conventions (long s, etc.).
  \item \textbf{Phoneme mapping with cognate-based ambiguity resolution:}
    Apply the default mapping table with overrides per cognate. The
    \textit{e} = /iː/ default (Goddard) is the most important rule.
    \textit{gh} is resolved per cognate (Schmick /x/ or /h/).
  \item \textbf{Morphological segmentation:} Identify and strip person
    prefixes (n-, k-, w-). Segment verbal suffixes (-uh, -eh, -seh, -oo,
    -nuh). Produce root forms for lexicon building.
  \item \textbf{Cross-validation:} Compare converted forms against Schmick's
    phonemicized forms, against Munsee cognates (Zeisberger, Goddard), and
    against PA reconstructions (Pentland, Siebert). Flag entries where
    converted phonemes mismatch expected reflexes.
\end{enumerate}

\subsubsection{Reference Cognate Set for Validation}

The following Edwards forms have well-understood cognates and should be used
as the validation baseline for the conversion:

\begin{table}[h]
\centering
\begin{tabular}{llll}
\toprule
\textbf{Edwards form} & \textbf{Gloss} & \textbf{PA reconstruction} & \textbf{Expected Mahican} \\
\midrule
\textit{neesoh} & `two' & *niːšwi & /niːsə/ \\
\textit{seepoo} & `river' & *siːpoːwi & /siːpə w/ \\
\textit{keesogh} & `sun' & *kiːšoʔθ a & /kiːsox/? \\
\textit{hpoon} & `winter' & *pepoon- & /pə puːn/? \\
\textit{hkeesque} & `eye' & & /kiːskw/? \\
\textit{nannah} & `five' & *nyaːlanwi & /naːnə w/? \\
\textit{penumpaufoo} & `boy' & & \\
\textit{pumiffoo} & `he walks' & *pemihθ eːwa & /pə mihseːwa/? \\
\textit{ncmeetseh} & `I eat' & *miːči- & /nə miːčə/? \\
\textit{ghusooh} & `eight' & *nešwaːšika & \\
\textit{noghkah} & `my father' & & /nə xk-/ \\
\textit{muhhekaneew} & `Mahican' & & /mə hiːkaniːw/ \\
\bottomrule
\end{tabular}
\caption{Reference cognate set for validation}
\label{tab:1369-cognates}
\end{table}

\subsubsection{Software Architecture}

Given the ambiguities and the need for comparative-evidence-based resolution,
the conversion system should be implemented as:

\begin{itemize}
  \item \textbf{FSM core} for the unambiguous mappings (consonants, clear
    vowels)
  \item \textbf{Cognate database} providing priors for ambiguous mappings
    (\textit{gh}, \textit{e} values, vowel length)
  \item \textbf{Resolver module} that queries the cognate database before
    applying ambiguous rules; falls back to a default with a confidence flag
    if no cognate found
  \item \textbf{Output format:} phonemic transcription with confidence
    markers: \texttt{[high]} = confirmed by cognate, \texttt{[med]} = default
    rule no contradicting cognate, \texttt{[low]} = guess no cognate
    available, \texttt{[FLAG]} = contradictory evidence manual review needed
\end{itemize}

\subsection{Conversion Workflow}
\label{sec:1369-workflow}

\texttt{[Pending implementation --- see Issue \#1369]}

The conversion workflow will implement the research approach described above
as a multi-stage pipeline. Each stage will be implemented as a separate
module with its own test suite.

\subsection{Linguist Feedback}
\label{sec:1369-feedback}

\texttt{[Template --- content pending review]}

The following aspects require linguist review:

\begin{itemize}
  \item \textit{gh} grapheme resolution: confirmation of /x/ vs.\ /h/
    mapping per cognate
  \item \textit{e} = /iː/ default: verification of the Goddard finding
    across all entries
  \item Morphological segmentation: review of prefix/suffix identification
  \item Cognate database: validation of reference cognate set
  \item Confidence markers: review of \texttt{[low]} and \texttt{[FLAG]}
    entries
\end{itemize}

\subsection{Related Work: Colonial Recorder Perception Problem}

The core challenge of this research card---English-speaking recorders
misperceiving and mis-transcribing Algonquian phonemes---is a known
phenomenon documented across Eastern North America. Two researchers' work
is directly relevant.

\textbf{Dr.\ Keith Cunningham} (Linguist, Nanticoke Indian Tribe). His
doctoral dissertation \emph{A Phonological Analysis of Nanticoke With
Practical Applications for Language Revitalization} (Georgetown University)
analyzes three 18th-century Nanticoke word lists (Heckewelder 1785) using
the same methodology: historical phonology from colonial recorders, with the
same challenges of multilingual informants and orthographic complexity. His
2024 Algonquian Conference presentation \emph{A Revised Phonological
Analysis of the Heckewelder Vocabulary of Nanticoke} directly parallels the
Edwards analysis---both involve 18th-century recorders with greater linguistic
awareness than their 17th-century predecessors, both document preaspiration
more reliably than earlier sources, and both require PA comparative
calibration. Cunningham's work demonstrates that the colonial recorder
perception problem extends beyond SNEA into the broader Eastern Algonquian
family.

\textbf{Dr.\ Craig Kopris} (Independent scholar). His paper \emph{Les sons
wyandots perçus par des oreilles étrangères} (Wyandot sounds perceived by
foreign ears, \emph{Recherches amérindiennes au Québec}, 2014) is the exact
framing of the Edwards problem: how European recorders systematically
misperceived and mis-transcribed indigenous phonemes. His \emph{Wyandot
Phonology: Recovering the Sound System of an Extinct Language} applies the
same recovery methodology to an unrelated language family (Iroquoian),
demonstrating that the foreign-ear perception problem is universal across
colonial documentation of North American languages. His Dictionary of
Wyandot project (NEH award FN-255576-17) shows the long-term application of
these methods to language revitalization.

The Edwards Mahican corpus (144 records) is unique among the SNEA sources in
being recorded by a fluent second-language speaker of the language (Edwards
learned Mahican as a child from playmates). This makes Edwards the most
phonetically reliable colonial recorder in the corpus---but the \textit{gh}
grapheme mystery and \textit{uh} reduction patterns documented in this card
show that even a fluent L2 speaker's orthography is filtered through English
phonological categories, the same problem that both Cunningham and Kopris
address in their respective language families.

\subsection{References}

\begin{itemize}
  \item Brasser, T.J.\ (1978). ``Mahican.'' In \emph{Handbook of North
    American Indians}, Vol.\ 15, \emph{Northeast}, pp.\ 198--212. Washington:
    Smithsonian Institution.
  \item Edwards, Jonathan. (1788). \emph{Observations on the Language of the
    Muhhekaneew Indians}. New Haven: Josiah Meigs. [Reprinted 1823 with notes
    by John Pickering, in Massachusetts Historical Collections.]
  \item Goddard, Ives. (1979). ``Munsee Delaware.'' In \emph{Handbook of
    North American Indians}, Vol.\ 15, pp.\ 228--234.
  \item Goddard, Ives. (1982). \emph{The Historical Phonology of Munsee.}
    International Journal of American Linguistics 48:16--48.
  \item Goddard, Ives. (2007--2008). Corrections to Pentland's Mahican
    historical phonology. Papers of the Algonquian Conference.
  \item Goddard, Ives. (2010). Handout on Mahican phoneme inventory.
    Conference on Language Reconstruction, Ann Arbor, MI.
  \item Masthay, Carl (ed.). (1991). \emph{Schmick's Mahican Dictionary,
    with a Mahican Historical Phonology by David Pentland}. Memoirs of the
    American Philosophical Society, Vol.\ 197. Philadelphia.
  \item Pentland, David. (1991). ``A Mahican Historical Phonology.'' In
    Masthay 1991.
  \item Warne, Janet. n.d. ``Time-Depth in Mahican Diachronic Phonology:
    Evidence from the Schmick Manuscript.'' McGill University.
  \item \emph{Important Mahican Vowel Information} (PDF). 2024.
    munseedelaware.com. Consultation with Carl Masthay and Ives Goddard.
  \item Zeisberger, David. (1772). \emph{Zeisberger's Delaware Vocabulary of
    1772}.
  \item Zeisberger, David. (1827). \emph{Grammar of the Language of the
    Lenni Lenape or Delaware Indians}. Philadelphia.
\end{itemize}
